\documentclass{report}
\usepackage[utf8]{inputenc}
\usepackage[T1]{fontenc}
\usepackage{geometry}
\usepackage{graphicx}
\usepackage{float}
\usepackage{hyperref}
\usepackage[english,indonesian]{babel}

\title{Optimasi Pemrosesan Template LaTeX Menggunakan Pendekatan Berbasis AI}
\author{Academic Support Agent}
\date{\today}

\begin{document}
\maketitle

\begin{abstract}

\end{abstract}

\chapter{PENDAHULUAN}


\chapter{TINJAUAN PUSTAKA}


% Original content (commented out):
% \documentclass[12pt]{report}%Author Kristiāns Kacars 2022. Last updated 08.02.2023
%\usepackage[utf8]{inputenc} neizmanto XeLatex, tad jāatkomentē.
\usepackage[T1]{fontenc}

\usepackage{geometry}
\geometry{
 a4paper,
 left=30mm,
 top=20mm,
 right=20mm,
 bottom=20mm,
 includehead
 }
 
\usepackage{xcolor}
\usepackage[]{graphicx}
\usepackage{setspace}
%\onehalfspacing %rindstarpas 1.5 vienības
\linespread{1.5}
\usepackage{placeins}%ērti izmantot \FloatBarrier, lai neļauti attēliem aiziet pārāk tālu no vajadzīgās vietas.
\usepackage{titlesec}%izveido atstarpes starp chapter un section nosaukumu
\titleformat{\chapter}
  {\normalfont\fontsize{16}{18}\bfseries\raggedright}{\thechapter.}{0pt}{\MakeUppercase}{}
  
\titleformat{\section}
  {\normalfont\fontsize{14}{16}\bfseries}{\thesection}{1em}{}
\titleformat{\subsection}
  {\normalfont\fontsize{12}{16}\bfseries}{\thesubsection}{1em}{}  

\titlespacing*{\chapter}{0pt}{0pt}{10pt}
\titlespacing*{\section}{0pt}{5pt}{5pt}


\usepackage{fancyhdr}%šis izveido header, lai tur ir līnija ar section nosaukumu
\newcommand{\changefont}{%
    \fontsize{9}{11}\selectfont
}
\pagestyle{fancy}
\fancyhf{}
\fancyfoot[C]{\thepage}
\fancyhead[L]{\changefont \rightmark}

\usepackage{polyglossia}
\setdefaultlanguage{latvian}
\usepackage{lipsum}
\setlength{\parskip}{1.5pt}%atstarpe starp rindkopām

%matematikas paketes
\usepackage{amsmath} 
\usepackage{amsthm}
\usepackage{amssymb}
\usepackage{bm}%ērts priekš treknraksta matemātiskajā vidē
\usepackage{esint}%integrāļu zīmes, kas nav parastajās ams paketēs

%--pašizveidotās matemātikas komandas
\newcommand{\bra}{\langle}
\newcommand{\ket}{\rangle}
%--

\newtheorem{theorem}{Teorēma}
\theoremstyle{definition}
\newtheorem{definition}{Definīcija}

\usepackage{siunitx}%mērvienībām

%Algoritmu atspoguļošanai
\usepackage[ruled, vlined, linesnumbered, algochapter]{algorithm2e}
\renewcommand{\algorithmcfname}{Algoritms}% nomaina nosaukumu uz latviešu variantu
\usepackage[labelfont={bf}]{caption}%attēlu un tabulu nosaukumu fonts

\usepackage[style=phys,%der arī citi, piem., chem-angew(bibliogrāfijā pie article nav links)
articletitle=false,
biblabel=brackets,%
chaptertitle=false,%
    sorting=none,%kārto pēc citēšanas secības
    date=short]{biblatex}
\addbibresource{biblio.bib}


\usepackage{hyperref}%Izveido spiežamus linkus saturam, bibliogrāfijām uc.
\hypersetup{
    colorlinks,
    citecolor=blue,
    filecolor=black,
    linkcolor=black,
    urlcolor=red
}


\newcommand*{\SavedEqref}{}%Dod krāsainus un spiežamus vienādojumu linkus
\let\SavedEqref\eqref
\renewcommand*{\eqref}[1]{%
  \begingroup
    \hypersetup{
      linkcolor=blue,
      linkbordercolor=black,
    }%
    \SavedEqref{#1}%
  \endgroup
}%vienādojumiem


\usepackage{chngcntr}%figure numeration
\counterwithin{figure}{chapter}
\counterwithin{table}{chapter}
\counterwithin{equation}{chapter}%vienadojumu numerācija veidojas pa nodaļām



%Pakete izveido apzīmējumu sarakstu
\usepackage{nomencl}
\makenomenclature
\renewcommand{\nomname}{Apzīmējumu saraksts}
%% Sadala grupās apzīmējumus
%Ja gribat pievienot papildu kategoriju, nokopējat-->
%   \ifstrequal{#1}{F}{Fizikas konstantes}{%
% tikai pēdējā rindiņā( kur ir Lielumu apzīmējumi) ieliekat papildu "}"
% -----------------------------------------
\usepackage{etoolbox}
\renewcommand\nomgroup[1]{%
  \item[\bfseries
  \ifstrequal{#1}{F}{Fizikas konstantes}{%
  \ifstrequal{#1}{N}{Skaitļu kopas}{%
  \ifstrequal{#1}{K}{Kaut kas cits}{%
  \ifstrequal{#1}{C}{Lielumu apzīmējumi}{}}}}%
]}
% --------------------------------------------
% ------- beidzas package izsaukšana un stila veidošana






%%%--------- Aizpildīt titullapu---------- 
%% Šis aizpildīs template titullapu ar jūsu informāciju
% Darba tips, i.e. bakalaura vai maģistra
\def\degree{Bakalaura darbs}
% Fakultātes nosaukums
\def\faculty{Fizikas, matemātikas un optometrijas fakultāte}
% Nodaļa iekš fakultātes
\def\department{Fizikas nodaļa}
% Universitātes nosaukums
\def\university{Latvijas Universitāte}
% Universitātes logo
\def\crest{\includegraphics[width = 0.5\textwidth]{LU_logo_LV_horiz.png}}%ja ieliekat attēlu kādā mapē, tad jāsauc attēls no šīs mapes, citādi neparādīsies tas titullapā.
\def\vietlaiks{Rīga, 2022}
\def\supervisor{Darba vadītājs: Dr. Jānis Bērziņš}
\def\studaplieciba{kk123456}
\author{Vārds Uzvārds}
\title{Bakalaura darba sagatave}
\date{Februāris 2022}
\newcommand{\thedate}[0]{21.03.2022}%iesniegšanas datums
%%%---------- Beidz aizpildīt titullapu----------- 




\begin{document}



\thispagestyle{empty}
\makeatletter
   \begin{center}
       \vspace*{1cm}
        
    \vspace{10mm}
    {\Large LATVIJAS UNIVERSITĀTE\\
    \MakeUppercase{\faculty}\\
    \vspace{2mm}
    \MakeUppercase{\department}}
    \vspace*{10mm}
    
    
    
    \vspace{5mm}
    {\Large \MakeUppercase{\textbf{\@title}}}
    \vspace{5mm}
    

       \vspace{1cm}
    \Large
    \MakeUppercase{\degree}
    \end{center}
    \vspace{3cm}
    \begin{flushleft}
    \large
       Autors: \textbf{\large \@author}\\
       Studentu apliecības Nr.: \studaplieciba \\
       \supervisor
    \end{flushleft}

       \vfill
     
    \begin{center}
    \Large      
    \MakeUppercase{\vietlaiks}
   \end{center}
\makeatother

\newpage




\thispagestyle{empty}
\noindent \textbf{Anotācija}

\noindent Anotāciju sagatavo latviešu un angļu valodā. Pēc saskaņošanas ar studiju
programmas direktoru var sagatavot arī papildu anotāciju kādā citā Eiropas Savienības oficiālajā
valodā. Anotācijā izklāsta problēmas būtību, pētījuma mērķus, uzdevumus, raksturo iegūtos
rezultātus. Anotācijas apjoms ir noteikts līdz 850 rakstu zīmēm, ieskaitot atstarpes.


\vspace{4mm}
\noindent \textbf{Atslēgas vārdi}: Atslēgvārdi raksturo darba tematu, rezultātus un izmantotās metodes. Ieteicamais
atslēgvārdu skaits ir no 4 līdz 8.

\vspace{20mm}
\noindent \textbf{Abstract}

\noindent In order to separate the equations, the radial part is set equal to a constant, and the form of the constant on the right above reflects the nature of the solution of the colatitude equation which yields the orbital quantum number. where $L_n^l$ is the associated Laguerre function.

\vspace{4mm}
\noindent \textbf{Keywords}: Thesis template, FMOF

\newpage
\tableofcontents
\newpage
%Apzīmējumu saraksts
\nomenclature[N]{\(\mathbb{R}\)}{Reālo skaitļu kopa}
\nomenclature[N]{\(\mathbb{H}\)}{Kvaternioni}

\nomenclature[F]{\(\hbar\)}{Planka konstante}
\nomenclature[C]{\(\psi\)}{Viļņa funkcija}

\printnomenclature



\addcontentsline{toc}{chapter}{Apzīmējumu saraksts}
\chapter*{Ievads}
\addcontentsline{toc}{chapter}{Ievads}
Ievadā pamato temata izvēli un aktualitāti, kā arī apraksta:

\begin{enumerate}
    \item pētāmās problēmas vai izvirzītās hipotēzes;
    \item darba mērķi un uzdevumus;
    \item izmantotās pētniecības metodes;
    \item faktoloģiskā materiāla avotus;
    \item darba struktūru.
\end{enumerate}



\begin{center}
    \chapter*{Lietošanas pamācība}
    \addcontentsline{toc}{chapter}{Lietošanas pamācība}
\end{center}
\section{Noformējums}
\noindent Kad vairs nevajag šo nodaļu, dzēšat nodaļu "Lietošanas pamācība" ārā.

Lai izveidotu titullapu, aizpildiet .tex faila sākumā doto tekstu "\textbf{Aizpildīt titullapu}".

Teksts tiek sadalīts nodaļās, sekcijās un apakšsekcijās. To svarīgums no augstākās ejot uz zemāko ir:
\begin{itemize}
    \item chapters, section, subsections
\end{itemize}
Ja vajadzīga nodaļa vai sadaļa, kura netiek nummurēta, bet to vajag saturā izmantot komandas
\begin{enumerate}
    \item \textbackslash chapter*\{Iedaļa\}
    \item \textbackslash addcontentsline\{toc\}
    \item \{chapter(iedaļas veids)\}\{Iedaļa\}.
\end{enumerate}

%Piemērs:
%    \chapter*{Lietošanas pamācība}
%    \addcontentsline{toc}{chapter}{Lietošanas pamācība}

\noindent Teksta stili ir
\begin{itemize}
    \item \textbf{treknrakstā} ar \textbackslash textbf\{..\}
    \item \textit{slīprakstā} ar \textbackslash textit\{..\}
    \item \textit{\textbf{treknā slīprakstā}} ar \textbackslash textit\{\textbackslash textbf\{..\}\}.
\end{itemize}



\noindent Lai pievienotu apzīmējumu sarakstam apzīmējumu izmanto
\begin{itemize}
    \item \textbackslash nomenclature[N]\{apzīmējums\}\{Paskaidrojums\},
\end{itemize}
kur [N] parāda kam pieder apzīmējums. Pagaidām ir \textbf{Lielumu apzīmējumi} ar [C], \textbf{Fizikas konstantes} ar [F] un \textbf{skaitļu kopas} [N], ja vajag jaunu kategoriju, pie \textbackslash renewcommand\textbackslash nomgroup pievieno jaunu kategoriju (neaizmirst pielikt \}, lai noslēgtu pievienoto kategoriju, citādi nekompilēsies).

Lai labāk sakārtots overleaf fails iesaku izmantot mapes, kur ir attēli salikti pa nodaļām, palīdzēs neapmaldīties.
\section{Matemātiskais pieraksts}



Matemātiskās funkcijas vai darbības rakstīt ar \textbackslash... kā $ \sin,\ \lim$ nevis $sin,\ lim$. Izmanto \textbackslash mathbb, lai iegūtu $\mathbb{R}$, $\mathbb{C}$. Vektorus var apzīmēt ar bultiņu vai treknrakstā: \[\vec{x}\quad \text{ar \textbackslash vec vai ar \textbackslash bm}\quad \bm{x},\] bet tas pie apzīmējumiem jāpaskaidro.

Vienādojumus var rakstīt "inline mode" - pašā rindiņā - ar \$\$  vai \textbackslash(... \textbackslash), piemēram, $y = x$ \( \lim_{n\xrightarrow{}\infty} x_n = x\), vai "display mode", kad vienādojums ir atseviķi no teksta
\begin{equation}
    \int\limits_{0}^{\infty}\frac{\sin x}{x}dx \quad \text{ja vajadzīgs komentārs piem.}\ \forall  x\in \mathbb{R}.
\end{equation}
Vienādojumus var arī līdzināt ar \textbackslash begin\{align\}
\begin{align}
	a &= 5 \\
	b + c &= a \\
	a -2\cdot3 &= 5/4.
\end{align}
Līdzīgi align strādā arī gather, tur pats tex līdzina vienādojumus. Ja vairākus vienādojumus grib zem vienas references kā vienādojumā \eqref{eq1}, tad izmanto aligned vai gathered, tie ir iekš begin\{equation\} vides:
\begin{equation}
    \begin{gathered}
    (x-3)x = 2x+1\\
    x^2-5x-1 = 0.
    \end{gathered}
    \label{eq1}
\end{equation}
Vēl vienādojumus display mode var rakstīt ar \textbackslash[...]\textbackslash: \[ \iiint\limits_V \nabla\cdot \bm{E}\,dV = \oiint\limits_S \bm{E}\cdot \, d\bm{A}.\]
Kad raksta vienādojumus ar iekavām izmantot \textbackslash left[.. \textbackslash right]:
\begin{equation}
    \left[\frac{x+y}{x-y}\right]\quad \text{nevis}\quad [\frac{x+y}{x-y}]
\end{equation}.

Matricu rakstība ir ar pmatrix vidi (tas pats tikai citas iekavas ap matricu ir arī bmatrix, Bmatrix, vmatrix, Vmatrix): 
\begin{equation}
\begin{pmatrix}
a_{11}&a_{12}&\dots&a_{1n}\\
a_{21}&a_{22}&\dots&a_{2n}\\
\hdotsfor[2]{4}\\
a_{n1}&a_{n2}&\dots&a_{nn}
\end{pmatrix}.
\end{equation}

Ja vajag figūriekavas, kur atšķir gadījumus izmanto \textbackslash begin\{cases\}:
\begin{equation}
    f(x)=\begin{cases}
0,& \text{ja $x>0$},\\
x, &
\text{ja $x<0$}.
\end{cases}
\end{equation}
Ja grib atsaukties uz kādu vienādojumu, tad tam ir jābūt \textbackslash label\{nosaukums\}, un uz to atsaucas izmantojot \textbackslash eqref\{nosaukums\}, tas dos šo \eqref{eq1}.


Ja vajag definēt jaunas komandas piemēram bra-ket notācijai, tad tex faila sākumā, pēc matemātikas paketēm izsaukšanas definē\newline \textbackslash newcommand\{\textbackslash ket\}\{\textbackslash langle\}, kas ļauj vieglāk uzrakstīt: $\bra\psi|\hat{H}|\psi\ket$. Kvantu mehānikā bieži izmanto īpašas konstantes un mērvienības: $\hbar$, \AA. 

Failā definētas komandas teorēmas un pierādījumi:
\begin{theorem}
Saistīti stāvokļi viendimensionālos potenciālos nav deģenerēti.
\end{theorem}
\begin{proof}
$\psi_1\neq \psi_2$
\end{proof}
Ir arī definīcijas, var ielikt arī lemmas un citas lietas, atkarībā no vajadzības.
\begin{definition}
Par Hilberta telpu sauc pilnu normētu telpu ar $||x|| = \sqrt{\bra x|x\ket}$
\end{definition}
\section{Noderīgi instrumenti}
\subsection{Atsauču veidošana}
Bibliogrāfiju veido sagatavojot .bib failu (skatīt biblio.bib).
Atsaucoties uz zinātniskajiem rakstiem, ir parocīgi izmantot \textit{Export citation as BibTex} rīku. Tas izveidos .bib failā ieliekamu atsauci. Atsaukties var arī uz grāmatām \cite{dirac}, mājaslapām \cite{Algorithm}, piemēru apskatīt .bib failā. Var arī atsaukties uz vairākiem avotiem vienlaikus \cite{Algorithm, AlgorithmLV}.


\subsection{Attēlu ievietošana}
\FloatBarrier
Attēlus ievieto izmantojot \textbackslash begin\{figure\}. Ieteicams izmantot vektorgrafika attēlus (.pdf nevis .png). Piemēram, uz to var atsaukties kā \textbackslash ref\{..\}: attēls \ref{fig:LU logo}.
\begin{figure}[!ht]
    \centering
    \includegraphics[width=0.5\textwidth]{LU_logo_LV_horiz.png}
    \caption{LU logo}
    \label{fig:LU logo}
\end{figure}
Ja vēlas 2 attēlus blakus izmanto minipage
\begin{figure}[!ht]
    \centering
    \begin{minipage}{0.5\textwidth}
        \centering
        \includegraphics[width=0.4\textwidth]{LU_logo_LV_horiz.png} % first figure itself
        \caption{Pirmais attēls}
        \label{fig:min1}
    \end{minipage}\hfill
    \begin{minipage}{0.5\textwidth}
        \centering
        \includegraphics[width=0.4\textwidth]{LU_logo_LV_horiz.png} % second figure itself
        \caption{Otrais attēls}
        \label{fig:min2}
    \end{minipage}
\end{figure}
\FloatBarrier
Atsaukties uz tiem var atseviški, katram vajag tikai savu \textbackslash label, attēls \ref{fig:min1} un \ref{fig:min2}. Minipage var ielikt arī gan vienu attēlu, gan otru tabulu, netikai vairākus attēlus.
\subsection{Tabulas un saraksti}

Tabulas var veidot ar \textbackslash begin\{table\} kā tabulā \ref{tab1}. Pārvietot \textbackslash caption virs tabulas, nevis zem kā ir automātiski, jo virsrakstu vajag virs tabulas.

\begin{table}[!ht]
\centering \caption{Pirmā tabula.}
\begin{tabular}{ |c|c|c| }% c automātiski padara kolonnu tik garu cik vajag
\hline
Npk. & Dati 1, \unit{kg.m.s^{-1}} &Dati 2, \unit{\mu W} \\
\hline
1. & 10.14 & 0.89\\
\hline
2. & 9.85   & 0.3 \\
\hline
3. &11.39 & 0.2 \\
\hline
\end{tabular}
\label{tab1}
\end{table}
\FloatBarrier
Sarakstus var veidot ar \textbackslash begin\{.....\}:
\begin{itemize}
    \item itemize (ar puntkiem)
    \item enumerate (ar skaitļiem),
    \item var taisīt arī sarakstu sarakstā
    \begin{itemize}
        \item[!--] cits saraksts 
        \item[!--] ar paša izdomātu numerācijas stilu
    \end{itemize} 
\end{itemize}

\subsection{Algoritmu rakstīšanas pakete}
Ja vajadzīgs uzrakstīt algoritma darbības principu, kā algoritmā \ref{alg:typeChecking}, var izmantot paketi algorithm2e. Vairāk par šo paketi var apskatīt \cite{Algorithm}, un ja vajag pārlikt kādu keyword uz latviešu valodu var noderēt \cite{AlgorithmLV}.

\begin{algorithm} [!ht]
	\caption{Mans algoritms} \label{alg:typeChecking}
	\KwIn{Polinoma pakāpe}
	\KwResult{Atgriež polinomu, ar prasīto pakāpi $n$.}
	\SetKwData{s}{s}
	\BlankLine
	\While{Kaut kas tiek darīts}{
		\lForEach{global statement \s} {
		}
		\ForEach{function $fun$} {
		\lForEach{statement \s in $fun$} {
		}
		}
		\lForEach{statement \s in game} {
		}
	}
	%\eIf{error messages were found}{\Return \False\;}{\Return \True\;}
\end{algorithm}
\begin{center}
    \chapter{Teorija}
   
\end{center}
\section{Pašsaistītā cikla metode}
No kvantu mehānikas ir zināms, ka pirmie pašsaistīto cikla metodi izmantoja Kons un Šems \cite{PhysRev.140.A1133}. Šī metode ar visiem izvedumiem aprakstīta \cite{dirac}. Sniegšu īsu ievadu metodes darbības principā, lai pamatotu iegūtos rezultātus.


\begin{center}
    \chapter{Nākamā nodaļa}
\end{center}
\section{Pirmā sadaļa}

\section{Otrā sadaļa}
\subsection{Pirmā apakšsadaļa}




\begin{center}
    \chapter{Rezultāti}
   
\end{center}
\section{Rezultātu analīze}
Rezultātu daļā parāda svarīgākos iegūtos rezultātus, salīdzina tos ar līdzīgiem
pētījumiem, novērtē rezultātu atbilstību izvirzītajām problēmām vai hipotēzei. 
\chapter*{Secinājumi}
\addcontentsline{toc}{chapter}{Secinājumi}
Noslēguma secinājumos atšķirībā no secinājumiem katras nodaļas nobeigumā (ja
autors tos formulējis) izsaka plašāku vispārinājumu, iesaka pētāmo problēmu risinājumu un
norāda turpmākos pētījuma virzienus. Apjomīgākos pētījumos secinājumus ieteicams strukturēt.

\cleardoublepage
\phantomsection
\addcontentsline{toc}{chapter}{Bibliogrāfija}
\printbibliography

\newpage
\chapter*{Pielikums A}
\addcontentsline{toc}{chapter}{Pielikums A}

Ja darbam nepieciešams, dažādus palīgmateriālus var ievietot pielikumā. Tajā
parasti iekļauj aprēķinu starprezultātus, ilustrācijas, anketu paraugus, kartes, aparātu un ierīču
aprakstus u. c.



\newpage

\thispagestyle{empty}

\begin{center}
     \textbf{\LARGE Pateicība}
\end{center}

Pateicību izsaka personām, kuras ir sniegušas organizatorisku vai finansiālu
palīdzību materiālu vākšanā vai arī metodiskus ieteikumus un morālo atbalstu darba izstrādē. Nav obligāts.

\newpage


\thispagestyle{empty}

\begin{center}
     \textbf{\LARGE Autora ieguldījums}
\end{center}

Nav obligāts.


%dokumentācijas lapa
% autoram/ei izskatīt, nomainīt locījumu vārdiem
\newpage
\thispagestyle{empty}
\makeatletter
{\setstretch{2.0}
\degree \ "\@title" izstrādāts LU Fizikas, matemātikas un optometrijas fakultātē.
    
    
\vspace{10mm}
Ar savu parakstu apliecinu, ka pētījums veikts pastāvīgi, izmantoti tikai tajā norādītie informācijas avoti un iesniegtā darba elektroniskā kopija atbilst izdrukai:

Autors: \rule{30mm}{0.2mm} \@author
    
\vspace{20mm}
Rekomendēju/nerekomendēju darbu aizstāvēšanai

Vadītājs/a: \supervisor \ \rule{30mm}{0.2mm} \thedate
   
\vspace{20mm}
\noindent Recenzents: \rule{60mm}{0.2mm}

\vspace{20mm}
Darbs iesniegts Fizikas nodaļā \rule{35mm}{0.2mm} \thedate

Dekāna pilnvarotā persona: \rule{60mm}{0.2mm}

\vspace{20mm}

Darbs aizstāvēts bakalaura gala pārbaudījuma komisijas sēdē

\rule{7mm}{0.2mm}.\rule{7mm}{0.2mm}.\rule{14mm}{0.2mm} .prot. Nr. \rule{7mm}{0.2mm}

Komisijas sekretāre: \rule{60mm}{0.2mm}

}
\makeatother








\end{document}


\end{document}